\item Para pruebas bacteriológicas de los suministros de agua y en clínicas médicas, las muestras se deben incubar rutinariamente durante 24 h a $37^\circ\text{C}$. El voluntario del Cuerpo de Paz e ingeniero del MIT, Amy Smith, inventó un incubador de bajo costo y bajo mantenimiento. El incubador consiste en una caja aislada con espuma que contiene un material ceroso que se funde a $37.0^\circ\text{C}$, intercalado entre tubos, platos o botellas que contienen las muestras de prueba y medio de cultivo (alimento para bacterias). Fuera de la caja, el material ceroso primero se funde mediante una estufa o colector solar. Luego el material ceroso se pone en la caja para mantener calientes las muestras de prueba, mientras se solidifica. El calor de fusión del cambio de fase del material es $205\text{ kJ/kg}$. Modele el aislamiento como un panel con $0.490\text{ m}^2$ de área superficial, 4.50 cm de espesor y $0.012\text{ W/m}\cdot^\circ\text{C}$ de conductividad. Suponga que la temperatura exterior es de $23.0^\circ\text{C}$ durante 12.0 h y $16.0^\circ\text{C}$ durante 12.0 h. (a) ¿Qué masa del material ceroso se requiere para realizar la prueba bacteriológica? (b) Explique por qué su cálculo puede efectuarse sin conocer la masa de las muestras de prueba o del aislante.
